\section{Pregunta de Investigación}
A partir del problema descrito, la pregunta que guía esta investigación es la
siguiente: ¿en qué medida un sistema de acceso con código QR dinámico, leído
por un ESP32-CAM, reduce el tiempo promedio de ingreso por estudiante
respecto a la revisión visual actual, y en qué medida disminuye los ingresos
con credenciales ajenas? Con ello se busca saber si el sistema ayuda a
aliviar los estragos del control de acceso actual del TESCHA. Para
responderla se medirán tres indicadores, antes del piloto y durante él, con
la hoja del Apéndice~C: el tiempo promedio de validación por estudiante, la
longitud máxima de la fila en el horario pico y el número de incidencias de
identidad detectadas.
